% =====================================================================
\section{Limiti e prospettive}
\label{sec:limiti}
% =====================================================================

\begin{itemize}
  \item \textbf{Adozione.} La 015 è usata in reti reali, la 018 quasi solo
        in testbed e prototipi~\cite{madqci2024}.
  \item \textbf{Multi-dominio incompleto.} Lo scenario normativo è
        QKD$+$OTN; il caso QKD$+$QKD richiede estensioni come quelle di
        Sim et al.\ e il coordinamento con la GS~QKD~020 per il passaggio
        delle chiavi tra KMS.
  \item \textbf{Metriche.} L'ESKR è un indicatore istantaneo e non tiene
        conto della chiave già accumulata nei buffer né della sua
        dinamica. Laschet et al.~\cite{laschet2026} propongono indicatori
        più ricchi per la selezione dei cammini.
  \item \textbf{Riservatezza della topologia.} Esporre la topologia
        interna all'orchestrator è accettabile dentro un operatore, meno
        tra operatori diversi; servono schemi di instradamento inter-dominio
        che preservino la riservatezza~\cite{laschet2026,interconnect2026}.
  \item \textbf{Sicurezza dell'interfaccia.} Lo standard non specifica i
        meccanismi di protezione del canale orchestrator–controller, che
        restano a carico dell'implementazione.
  \item \textbf{Algoritmi.} Il calcolo dei cammini non è standardizzato:
        è un vantaggio per l'innovazione~\cite{routing2024,mendez2026}, ma
        rende meno prevedibile il comportamento in scenari multi-vendor.
\end{itemize}

% =====================================================================
\section{Conclusioni}
\label{sec:conclusioni}
% =====================================================================

Letta in sequenza, la produzione ETSI sul controllo della QKD ha una
logica chiara. La GS~QKD~002 ha mostrato che la QKD utile è una rete, con
gestione delle chiavi, requisiti di servizio differenziati e integrazione
nella rete ottica dell'operatore. La GS~QKD~015 ha risposto con
l'astrazione del nodo SD-QKD: moduli come interfacce, un key manager unico,
link fisici e virtuali, applicazioni con QoS, un controller che registra,
pianifica e configura, tutto descritto in YANG. La GS~QKD~018 ha aggiunto
il livello di orchestrazione: ogni dominio diventa un grafo astratto, le
risorse di chiave vengono monitorate tramite l'inventario e il provisioning
segue uno schema in cui il controller propone e l'orchestrator decide. Il
disegno ricalca ACTN, riusa YANG, NETCONF e RESTCONF e mantiene il principio
di sicurezza essenziale: le chiavi non lasciano mai i nodi. Le esperienze di
KOREN, del Politecnico di Milano e di MadQCI mostrano che lo schema funziona
ma che l'orchestrazione QKD è ancora giovane. Il passo successivo è
estenderla davvero al multi-operatore, con metriche migliori e garanzie di
riservatezza tra domini.
